\head{第十三版新核查：第0章22.1--22.3}
\source{EGA IV第一部分，分册印刷第182--193页的相关部分／上传PDF第181--192页；作者正式勘误见第四卷第四部分原分册第345--348页的相关条目}
本轮从用户已收到的第十二版工程接续。将该工程的各个输入文件展开后，所得源码与已交付第十二版单文件逐字相同；其真实接续点为§22章首及（22.1.1）。本轮连续译入22.1--22.3，止于（22.3.4）及证明。22.4开始的内容虽已用于辨认下一接续位置，但不计入本版译文。

新增正文采取校订表述，原文差异、修改理由与出处放在脚注；本报告给出必要的反例、补证和适用范围核查。编号M13为本轮追踪编号。作者正式修订、译者对独立陈述的范围修正、上下文澄清及补充证明分别列出，不能按脚注数统计“错误定理”。历史报告保持原文，未因本轮合订而重新获得数学审定。

\section*{核查索引}
\addcontentsline{toc}{section}{核查索引}
\begin{longtable}{@{}p{.17\linewidth}p{.19\linewidth}p{.55\linewidth}@{}}
\toprule 标识&原条目&本轮处理\\ \midrule\endhead
M13-E01&22.2.6&分类唯一性限于完备环。给出非完备情形的反例；保留原文前后语境的说明。\\
M13-N03、N05&22.2.3、22.2.8&“替换条件（ii）”仍须保留（i）；给出脱离（i）时的检验例。这是作用域澄清。\\
M13-G03&22.1.4&展开阿廷纤维积与平方零核的递归提升，不借用有问题的旧引理表述。\\
M13-A01&22.1.1&依作者正式勘误改引20.6.10。\\
M13-A02&22.2.7&依作者正式勘误补明预进制拓扑。\\
M13-A03&22.3.3&依作者正式勘误将 $\widehat B_{\mathfrak p}$ 改为 $\widehat B_{\mathfrak q}$。\\
M13-T01&22.2.10&明确“$k$-扩张、相对于 $k_0$ 平凡”与“$k_0$-扩张”的差别。\\
其余M13-G、N项&相应脚注&展开线性对偶、标量结构、同构类唯一性、局部化幂次条件及非循环依赖；不当作反例数量。\\
\bottomrule
\end{longtable}

\audit{M13-E01}{分类定理须限定为完备局部代数}{原独立陈述缺少完备性；前文语境已指明完备情形}
\oldauditsrc{$0_{\mathrm{IV}}$（22.2.6），原分册第188页，并接第189页首句}
原定理先要求 $A$ 为“局部诺特 $k$-代数”，没有在该句中写“完备”，随后断言它由维数决定，允许的同构须在剩余域上诱导恒等。但前文（22.2.4）开头明确讨论诺特完备局部代数，（22.2.5）也只分类这样的对象。因此不能把省去“完备”的单独陈述照字面推广到所有局部诺特环。

\lead{反例。}
取有限域 $k=\mathbf F_p$，令
\[
 A_1=k[t]_{(t)},\qquad A_2=k[[t]],\qquad K=k.
\]
二者都是一维正则诺特局部环，极大理想由 $t$ 生成，剩余域都是 $k$。而且
\[
 \Upsilon_{K/k}=0,
\]
所以原陈述的非完美性模有限性条件也成立。

逐项核对其形式光滑性如下。对任意离散 $k$-代数 $R$、平方零理想 $I$ 和连续 $k$-同态 $A_i\to R/I$，$t$ 的像 $\bar b$ 为幂零元：连续性使某个 $(t)^N$ 落入核。任选提升 $b\in R$，则 $b^N\in I$，所以 $b^{2N}=0$。

对于 $A_1$，以 $t\mapsto b$ 定义多项式环同态。凡分母 $g(t)$ 满足 $g(0)\ne0$，其值 $g(b)$ 为非零常数加幂零元，因而可逆，所以同态延拓到 $k[t]_{(t)}$。它在 $(t)^{2N}$ 上为零，故连续，并提升原同态。

对于 $A_2$，用有限和
\[
 \sum_{r\ge0}a_rt^r\longmapsto\sum_{r=0}^{2N-1}a_rb^r
\]
定义所需同态；由于 $b^{2N}=0$，乘法相容且连续，模 $I$ 后为原映射。因此二者均对各自预进制拓扑形式光滑。

然而 $A_1$ 可数：多项式集合可数，分式可由两个多项式编码。$A_2$ 不可数：二元序列 $(a_r)_{r\ge0}$ 通过 $\sum a_rt^r$ 单射到它，序列集合由Cantor对角论证不可数。故 $A_1$ 与 $A_2$ 甚至不可能存在集合双射，更不可能有环同构。

\lead{校订后的范围。}
本版将（22.2.6）写在\emph{诺特完备局部} $k$-代数类中，沿用其原证明。若仍允许非完备 $A$，存在性判据本身不变，但分类唯一性应针对 $\widehat A$：由形式光滑性与完备化的相容性，$\widehat A$ 落入（22.2.5）的范围，而不同的原局部环可以有同一完备化。

本例不反驳（22.2.5），也不是声称原作者意图分类非完备环。它说明在供独立引用的正文命题中，必须把上下文的完备性限定写出。

\audit{M13-N03}{替换条件时不能丢掉原定理的另一个条件}{作用域提醒；不把正确上下文误判为错误}
\oldauditsrc{（22.2.3），第186页；（22.2.8），第189--190页}
原定理（22.2.2）的形式是
\[
 A/k\text{形式光滑}
 \quad\Longleftrightarrow\quad
 \bigl[(\mathrm i)\ \text{相伴分次对称代数同构}\bigr]
 \ \land\
 \bigl[(\mathrm {ii})\ \chi_{A/k}\text{单射}\bigr].
\]
随后的“条件（ii）可以替换为……”是在保留（i）的前提下说的。原证明也写明a)与（i）的合取等价于d)与（i）的合取。它并未证明脱离（i）后a)、b)、c)、d)仍全部等价。

\lead{一个能区分两种读法的实例。}
令
\[
 k=\mathbf F_2(s),\qquad
 A=k[x,y]/(x^2-sy^2,\ y^3),\qquad
 \mathfrak m=(x,y).
\]
由于 $x^4=s^2y^4=0$ 且 $y^3=0$，$\mathfrak m$ 是唯一素理想，$A$ 为阿廷局部环，剩余域为 $k$。把它写成
\[
 A=\bigl(k[y]/(y^3)\bigr)[x]/(x^2-sy^2)
\]
可见 $A$ 在 $k[y]/(y^3)$ 上自由，基为 $1,x$，因此
\[
 1,y,y^2,x,xy,xy^2
\]
为 $k$-基，特别是 $y^2\ne0$。又有 $\mathfrak m^4=0$。

因为 $K=A/\mathfrak m=k$，$\Upsilon_{K/k}=0$，所以（22.2.2，ii）自动成立。条件a)也成立：增广 $A\to k$ 在模掉 $\mathfrak m$ 后给 $\Omega^1_k\to\Omega^1_A\otimes_A k$ 一个左逆。

可是 $\Omega^1_k=k\,ds$，对关系 $x^2-sy^2=0$ 作绝对微分，在特征2中得到
\[
 y^2\,ds=0\quad\text{于 }\Omega^1_A.
\]
而 $y^2ds\ne0$ 于 $\Omega^1_k\otimes_k A=A\,ds$。故
\[
 \Omega^1_k\otimes_k A\longrightarrow\Omega^1_A
\]
不单射，尤其不可能有左逆。取 $n=3$ 时 $A/\mathfrak m^{n+1}=A$，所以（22.2.3，b）失败，c)也失败。

这里（i）确实不成立：$\mathfrak m/\mathfrak m^2$ 为二维，而对称代数有任意高次非零部分，$\gr_{\mathfrak m}(A)$ 却在次数至少4时为零。由此可见，问题在于删掉（i）后的误读，而不是原定理结论错误。

\lead{本版处理。}
（22.2.3）和（22.2.8）都直接写明“在保留条件（i）时”。上述检验例仅收入报告；脚注告知它的作用，不用它推翻原文在正确作用域内的陈述。

\audit{M13-G03}{阿廷测试环如何给出相容的逐级提升}{补足原证明的构造；结论不变}
\oldauditsrc{（22.1.4），第185--186页}
原证明已将问题化为：两个完备诺特局部 $A$-代数 $B,B'$，已给定闭纤维同构 $v_0:B/\mathfrak mB\to B'/\mathfrak mB'$；须构造同时与上一阶映射及该闭纤维同构相容的映射
\[
 v_j:B/\mathfrak n^{j+1}\longrightarrow B'/\mathfrak n'^{\,j+1}.
\]
若只分别选取各阶提升，不能保证相容，因而不能直接取逆极限。

固定 $j\ge1$，令 $R_j=B'/\mathfrak n'^{\,j+1}$。原证明所需的两个商分别由 $B'$ 的理想
\[
 I_j=\mathfrak n'^{\,j},\qquad
 J_j=\mathfrak n'^{\,j+1}+\mathfrak mB'
\]
给出。将上一阶数据和闭纤维数据合成同态到
\[
 Q_j=(B'/I_j)\mathop{\times}_{B'/(I_j+J_j)}(B'/J_j).
\]
环同态 $B'\to Q_j$ 满射：给定 $(a\bmod I_j,b\bmod J_j)$，有 $a-b=i+j'$，其中 $i\in I_j$、$j'\in J_j$；于是 $a-i=b+j'$ 同时给出两个分量。其核为 $I_j\cap J_j$。所以 $R_j\to Q_j$ 满射，核为
\[
 N_j=(I_j\cap J_j)/\mathfrak n'^{\,j+1}.
\]
因 $N_j\subseteq\mathfrak n'^{\,j}/\mathfrak n'^{\,j+1}$，且 $2j\ge j+1$，有 $N_j^2=0$。

$Q_j$ 是阿廷局部环，剩余域仍为 $K$；假设中的提升性质因而适用于 $B\to Q_j$ 和 $R_j\to Q_j$。得到的局部同态 $B\to R_j$ 杀掉 $\mathfrak n^{j+1}$，故诱导 $v_j$，并同时满足两项要求。取逆极限得到 $v:B\to B'$。

$B'$ 对 $A$ 平坦，$v$ 模去 $\mathfrak m$ 为同构。再核对（19.7.1.4）所需的分离完备性：$\mathfrak mB\subseteq\mathfrak n$，所以 $\mathfrak m$-进制柯西列也是 $\mathfrak n$-进制柯西列；$B$ 的 $\mathfrak n$-进制完备性给出极限，而每个 $\mathfrak m^rB$ 闭，保证该极限具有所需的 $\mathfrak m^r$ 同余。分离性由 $\bigcap_r\mathfrak n^r=0$ 给出。相同论证适用于 $B'$。因此可用该引理断定 $v$ 为同构。

这个构造只使用本命题的阿廷提升假设与（19.3.10.2）的纤维积事实，没有调用（19.3.10.1）未经修订的过弱条件，也不需要处理旧报告H8-P01的更一般拓扑范围。

\audit{M13-T01}{特征同态所对应的结构域}{用定义域消除证明末句的结构混淆}
\oldauditsrc{（22.2.10），第190页末句}
这里的对象具有两层结构：它是 $k$-扩张，但在忘掉部分结构后作为 $k_0$-扩张平凡。因此相对特征同态的定义域为 $\Upsilon_{K/k/k_0}$。它可能非零，并不等于以 $\Upsilon_{K/k_0}$ 为定义域的绝对特征同态；后者因 $k_0$-平凡而为零。

例如取 $k_0=\mathbf F_p$、$k=\mathbf F_p(s)$、$K=\mathbf F_p(a)$，规定 $s=a^p$。在环 $E=K[\epsilon]/(\epsilon^2)$ 上赋予 $k$-代数结构
\[
 s\longmapsto a^p+\epsilon.
\]
这定义了域 $k$ 的同态，因为任何非零有理函数的分母模 $\epsilon$ 后都非零，因而在 $E$ 中可逆。该扩张在 $k_0$ 上由 $K\hookrightarrow E$ 分裂；但它的相对特征同态把 $ds\otimes1$ 映为 $\epsilon$，并不为零。正文末句改为“$k_0$-平凡 $k$-扩张的特征同态”，与该条陈述、定义域及自然性完全一致。

\audit{M13-A01--A03}{本轮落实的作者正式勘误}{三项已有作者修订，不作为新发现}
\oldauditsrc{第四卷第四部分，《勘误与增补（清单三）》A部第345--346页及B部第15项（第348页）}
\begin{enumerate}[label=\arabic*.,leftmargin=2.5em]
\item （22.1.1）第183页倒数第12行，将（20.6.20）改为（20.6.10）。这里讨论扩张的特征同态，不是域扩张微分模的消失。
\item （22.2.7）第189页第5行，把“形式光滑”改为“对预进制拓扑形式光滑”。离散拓扑与预进制拓扑下的提升问题不能混用。
\item （22.3.3）第192页倒数第13行，把误排的 $\widehat B_{\mathfrak p}$ 改为 $\widehat B_{\mathfrak q}$。这里改的是所论局部环的下标，原有基环 $B_{\mathfrak r}$ 不变。正确结构同态是
\[
 B_{\mathfrak r}\longrightarrow\widehat B_{\mathfrak q},
 \qquad\mathfrak r=\mathfrak q\cap B.
\]
\end{enumerate}
这些修改已经写入正文，原写法与勘误出处放入脚注，书末还补入相应的正式勘误译文。没有把只读取过的整份勘误清单算作全部译完。

\audit{M13-G07--G08}{局部化、完备化与证明依赖}{核对附加条件和非循环证明顺序}
\oldauditsrc{（22.3.1）--（22.3.4），第191--193页}
（22.3.1）只给\emph{相对}形式光滑性；（22.3.2）用正则性及子域族条件将其提升为所需的形式光滑性。对于正特征的局部化，幂次条件
\[
 A^{p^{h(\alpha)}}\subseteq A_\alpha
\]
不可在叙述中遗漏。设 $\mathfrak s_\alpha=\mathfrak s\cap A_\alpha$，则对 $t\notin\mathfrak s$，令 $q=p^{h(\alpha)}$，有
\[
 t^q\in A_\alpha\setminus\mathfrak s_\alpha,
 \qquad t^{-1}=\frac{t^{q-1}}{t^q}.
\]
因此 $A_{\mathfrak s}=(A_\alpha\setminus\mathfrak s_\alpha)^{-1}A$，有限性和幂次条件都可传递，分式域不变。这才完成原证明中的局部化归约。

（22.3.3）把结论归结到闭纤维，再调用紧随其后的（22.3.4）。后一条并不反过来使用（22.3.3），而是利用Cohen结构定理把完备局部整环表示为正则幂级数子环的有限代数，然后直接应用（22.3.2）。因此证明依赖为
\[
 \text{Cohen结构及（21.8.8）}
 \ \to\ (22.3.2)\ \to\ (22.3.4)\ \to\ (22.3.3),
\]
其中（22.3.2）还使用（22.3.1）及22.2判据。此处“推论”在编号上后置，不构成逻辑循环。

\section*{核查边界与接续位置}
\addcontentsline{toc}{section}{核查边界与接续位置}
本轮连续完成§22章首及22.1--22.3，包含20个原著编号条目（含22.2.5.2）。下一条为22.4章首及（22.4.1），位于上传PDF第192页、原分册印刷第193页。第0章22.4--22.7、§23、第IV章主体及第V卷预备笔记余文均未计作完成。

本报告没有宣称对已有全部译文重新逐句审定；H8-P01等历史待核事项继续保留。程序检验与版面检查只能确认其实际检查项目，不能替代独立数学审定。具体输出页数、编译与渲染结果以本次生成的\texttt{checks/CHECKS.md}为准。
